% Figure: the dual (Sabathe) combustion cycle on the p-v plane, with heat added
% partly at constant volume (2 to x) and partly at constant pressure (x to 3),
% shown against the pure Otto (all constant volume) and pure Diesel (all
% constant pressure) limits it interpolates between. Curves are schematic
% adiabats p v^gamma = const drawn between fixed volume stations.
\begin{figure}[htbp]
\centering
\begin{tikzpicture}
\begin{axis}[
  width=12cm, height=8.4cm,
  xlabel={specific volume $v$ (arb.)},
  ylabel={pressure $p$ (arb.)},
  xmin=0.6, xmax=9.5, ymin=0, ymax=70,
  axis lines=left,
  xtick=\empty, ytick=\empty,
  legend pos=north east, legend cell align=left,
  legend style={font=\scriptsize},
  tick label style={font=\scriptsize}, label style={font=\small},
  clip=false,
]
% compression ratio r = 8 : v2 = 1, v1 = 8. Adiabat p = C / v^1.4.
% Choose C so p1 = 1 at v1 = 8 : C = 1 * 8^1.4 = 18.38.
% p2 at v2 = 1 : p2 = 18.38.
% Constant-volume heat add 2 -> x raises pressure from 18.38 to 55 at v = 1.
% Constant-pressure heat add x -> 3 at p = 55 from v = 1 to v = 1.6 (cut-off).
% Expansion adiabat 3 -> 4 : p = C3 / v^1.4, C3 = 55 * 1.6^1.4 = 55*1.90 = 104.6.
% p4 at v = 8 : 104.6 / 8^1.4 = 104.6/18.38 = 5.69.
% Constant-volume rejection 4 -> 1 at v = 8 from 5.69 to 1.
\addplot[engblue, very thick, samples=80, domain=1:8] {18.38/(x^1.4)};
\addlegendentry{compression \& expansion adiabats}
\addplot[engblue, very thick, samples=80, domain=1.6:8] {104.6/(x^1.4)};
% constant-volume heat addition 2 -> x at v = 1
\draw[engred, very thick] (axis cs:1,18.38) -- (axis cs:1,55);
% constant-pressure heat addition x -> 3 at p = 55
\draw[engred, very thick] (axis cs:1,55) -- (axis cs:1.6,55);
% constant-volume heat rejection 4 -> 1 at v = 8
\draw[engamber, very thick, dashed] (axis cs:8,5.69) -- (axis cs:8,1);
% state markers
\addplot[engblack, only marks, mark=*, mark size=1.4pt] coordinates
  {(8,1) (1,18.38) (1,55) (1.6,55) (8,5.69)};
\node[font=\scriptsize, anchor=north east] at (axis cs:8,1) {1};
\node[font=\scriptsize, anchor=east] at (axis cs:1,18.38) {2};
\node[font=\scriptsize, anchor=south east] at (axis cs:1,55) {x};
\node[font=\scriptsize, anchor=south west] at (axis cs:1.6,55) {3};
\node[font=\scriptsize, anchor=west] at (axis cs:8,5.69) {4};
% annotate the two heat-addition legs
\node[engred, font=\scriptsize, anchor=west] at (axis cs:1.15,40)
  {const-$v$ burn};
\node[engred, font=\scriptsize, anchor=south] at (axis cs:1.3,56.5)
  {const-$p$ burn};
\end{axis}
\end{tikzpicture}
\caption[Dual (Sabath\'e) cycle on the pressure-volume plane]{The dual, or
Sabath\'e, cycle splits its heat addition into a constant-volume rise from state
2 to state x, followed by a constant-pressure stretch from x to state 3, before
the adiabatic expansion to state 4 and the constant-volume rejection back to
state 1. The pure Otto cycle is the limit in which all the heat enters along the
vertical constant-volume leg, and the pure Diesel cycle the limit in which all of
it enters along the horizontal constant-pressure leg; the dual cycle interpolates
between them and models a real compression-ignition engine more faithfully than
either, since combustion in such an engine begins nearly at constant volume as
the premixed charge burns and finishes at nearly constant pressure as the rest
injects. The pressure-rise ratio at state x and the cut-off ratio at state 3 are
the two parameters that place the cycle along the Otto-to-Diesel span.}
\label{fig:vol04ch03:dual-cycle-pv}
\end{figure}
